%
% 11-reell.tex -- Erweiterung von Q durch die reellen Zahlen
%
% (c) 2026 Prof Dr Andreas Müller
%
\subsection{Erweiterung von $\mathbb{Q}$ durch die reellen Zahlen}
Schon im Altertum war bekannt, dass die Quadratwurzel $\!\sqrt{2}$ von 
$2$ nicht rational ist.
Diese Entdeckung, die normalerweise den Pythatgoräern zugeschrieben wird,
\index{Pythatgoraer@Pythagoräer}%
widersprach der verbreiteten Ansicht, dass sich alle Zahlen durch
Brüche ausdrücken lassen.
In diesem Abschnitt werden zwei Konstruktionen betrachtet, mit denen
sich eine Erweiterung von $\mathbb{Q}$ finden lässt, die $\!\sqrt{2}$
enthält.

%
% Die reellen Zahlen
%
\subsubsection{Die reellen Zahlen}
Die Quadratwurzel $\!\sqrt{2}$ ist bei weitem nicht die einzige irrationale
Zahl.
Jede ganze Zahl, deren Primfaktorzerlegung einen Primfaktor in ungerader
Anzahl enthält, hat eine irrationale Quadratwurzel.
Auch viele kubische Wurzeln, die eulersche Zahl $e$ oder die
Kreiszahl $\pi$ sind irrational, es lassen sich aber immer gute
Approximationen durch Brüche geben.
\index{Kreiszahl}%
\index{pi@$\pi$}%
Leonhard Euler hat beretis 1737 bewiesen, dass $e$ irrational ist.
\index{Euler, Leonhard}%
Für $\pi$ ist dies erst Johann Heinrich Lambert 1761 gelungen.
\index{Lambert, Johann Heinrich}%
Diese Beispiele deuten darauf hin, dass es eine grosse Menge von
Zahlen gibt, die beliebig genau durch Brüche approximiert werden
können, aber trotzdem keine Brüche sind.

Das Problem kann gelöst werden, indem man eine Menge $\mathbb{R}$ von
Zahlen definiert, die beliebig genau durch Brüche approximiert werden
können.
Dabei muss aber auch berücksichtigt werden, dass eine solche Zahl
verschiedene Approximationen haben kann.
Zum Beispiel kann $\pi$ sowohl durch Dezimalbrüche
\[
3
, 3.1
, 3.14
, 3.141
, 3.1415
, 3.14159
, 3.141592
, 3.1415926
, 3.14159265
%, 3.141592659
%, 3.1415926592
%, 3.14159265926
\to
\pi
\]
als auch durch die Kettenbruchentwicklung
\index{Kettenbruch}%
\[
\frac{3}{1}
\approx
3.{\color{gray}00},
\frac{22}{7} \approx 3.14{\color{gray}28},
\frac{333}{106}\approx 3.1415{\color{gray}09},
\frac{355}{113}\approx 3.141592{\color{gray}92}
\to
\pi
\]
approximiert werden.
Es reicht daher nicht, eine Approximation zu haben.
Es wird ausserdem ein Kriterium benötigt, mit welchem man entscheiden
kann, ob zwei Approximationen die gleiche Zahl approximieren.

\begin{definition}[Cauchy-Folge]
Eine Folge $(a_n)_{n\in\mathbb{N}}$ von Zahlen $a_n\in\mathbb{Q}$ heisst
\emph{Cauchy-Folge}, wenn für jedes $\varepsilon>0$ ein $N\in\mathbb{N}$
\index{Cauchy-Folge}%
gibt derart, dass
\[
|a_n-a_m| < \varepsilon
\]
für alle $m,n> N$ gilt.
\end{definition}

Die Dezimalbruchapproximation von $\pi$ ist eine Cauchyfolge.
Sei $a_n$ der endliche Dezimalbruch mit $n$ Nachkommastellen, die
mit der Dezimalbruchentwicklung übereinstimmen.
Dann ist
$a_n-a_m$ ein Dezimalbruch, dessen erste $\min(n,m)$ Nachkommastellen
$0$ sind.
Insbesondere ist
\[
|a_n-a_m| < 10^{-\min(n,m)}.
\]
Sei also $\varepsilon>0$ gegeben und $N\in\mathbb{N}$ derart,
dass $10^{-N}<\varepsilon$.
Falls $n,m>N$ sind, folgt
\[
|a_n-a_m|
<
10^{-\min(n,m)}
<
10^{-N}
<
\varepsilon.
\]
Damit ist gezeigt, dass $(a_n)_{n\in\mathbb{N}}$ eine Cauchy-Folge ist.

Endliche und unendliche Dezimalbrüche konvergieren gegen rationale Zahlen
im Sinne der folgenden Definition.

\begin{definition}[konvergente Folge]
\index{konvergente Folge}%
Eine Cauchy-Folge $(a_n)_{n\in\mathbb{N}}$ konvergiert gegen $a\in\mathbb{Q}$,
auch geschrieben als
\[
\lim_{n\to\infty} a_n = a
\qquad\text{oder}\qquad
a_n\to a,
\]
wenn für jedes $\varepsilon>0$ ein $N\in\mathbb{N}$ existiert derart, dass
\[
|a_n-a|<\varepsilon
\]
für $n>N$.
\end{definition}

Nicht jede Cauchy-Folge konvergiert gegen eine rationale Zahl.
Die Approximationsfolge von $\!\sqrt{2}$ ist eine Cauchy-Folge, die nicht
gegen eine rationale Zahl konvergiert.

Mit Cauchy-Folgen kann man auch Rechnen.
In der Analysis wird der folgende Satz bewiesen.

\begin{satz}
Seien $(a_n)_{n\in\mathbb{N}}$ und $(b_n)_{n\in\mathbb{N}}$ Cauchy-Folgen.
Dann sind $(a_n+b_n)_{n\in\mathbb{N}}$, $(a_n-b_n)_{n\in\mathbb{N}}$ 
und $(a_nb_n)_{n\in\mathbb{N}}$ Cauchy-Folgen.
Falls $a_n\to a$ und $b_n\to b$, dann gilt auch
$a_n\pm b_n\to a\pm b$ und $a_nb_n\to ab$.
Falls $b\ne 0$ ist, ist auch $(a_n/b_n)_{n\in\mathbb{N}}$ ein Cauchy-Folge
und $a_n/b_n \to a/b$.
\end{satz}

\begin{definition}[Nullfolge]
Eine Cauchy-Folge $(a_n)_{n\in\mathbb{N}}$ heisst \emph{Nullfolge}, wenn
sie gegen $0$ konvergiert.
\index{Nullfolge}%
\end{definition}

Die Rechenoperationen und der Begriff einer Nullfolge ermöglichen jetzt,
Cauchy-Folgen miteinander zu identifizieren, die den gleichen Grenzwert
haben, selbst wenn dieser nicht eine rationale Zahl ist.

\begin{definition}[Äquivalenz von Cauchy-Folgen]
Zwei Cauchy-Folgen
$(a_n)_{n\in\mathbb{N}}$
und
$(b_n)_{n\in\mathbb{N}}$
heissen \emph{äquivalent}, in Zeichen $(a_n)\sim (b_n)$ wenn
$(a_n-b_n)_{n\in\mathbb{N}}$ eine Nullfolge ist.
\index{aquivalente@äquivalente Cauchy-Folgen}%
\end{definition}

Wenn zwei Cauchy-Folgen äquivalent sind, dann approximieren sie 
die gleiche Zahl.
Die gesuchte Erweiterung der Menge der rationalen Zahlen $\mathbb{Q}$
besteht also aus Cauchy-Folgen von rationalen Zahlen, wobei jedoch
äquivalente Folgen die gleiche Zahl repräsentieren.

\begin{definition}[Äquivalenzklassen]
\index{Aquivalenzklasse@Äquivalenzklasse}%
Eine \emph{Äquivalenzklasse} $X$ der Äquivalenzrelation $\sim$ von
\index{Aquivalenzrelation@Äquivalenzrelation}%
Cauchy-Folgen ist eine Menge von äquivalenten Cauchy-Folgen.
\end{definition}

Zwei Cauchy-Folgen
$(a_n)_{n\in\mathbb{N}}\in X$ 
und
$(b_n)_{n\in\mathbb{N}}\in X$ 
in einer Äquivalenzklasse $X$ von rationalen Cauchy-Folgen unterscheiden
sich nur durch eine Nullfolge, es gilt also
\[
\lim_{n\to\infty} (a_n-b_n) = 0.
\]
Die reellen Zahlen können jetzt als Äquivalenzklassen von Cauchy-Folgen
definiert werden.

\begin{definition}[Reelle Zahlen]
Die Menge der reellen Zahlen $\mathbb{R}$ ist die Menge der Äquivalenzklassen
von rationalen Cauchy-Folgen bezüglich der Äquivalenzrelation $\sim$.
\end{definition}

Diese Konstruktion zeigt, dass man mit reellen Zahlen rechnen kann, indem
man geeignete Approximationen dafür verwendet.
Welche Approximation verwendet wird, ist egal.
Wenn die Approximationsfolgen äquivalent sind, unterscheiden sich auch
Resultate von Rechnungen mit den Folgen durch Nullfolgen.
Die reellen Zahlen enthalten nach Konstruktionen alle Quadratwurzeln
und die Zahlen $e$ und $\pi$.
Jede Gleichung, für deren Lösung eine Approximationsfolge existiert,
kann daher in $\mathbb{R}$ auch gelöst werden.
So kann man zum Beispiel die Lösung der Gleichung
\[
y = 1 + \frac1y
\]
dadurch approximieren, dass man die Folge $(y_n)_{n\in\mathbb{N}}$
durch
\[
y_0 = 1,
\quad
y_{n+1} = 1 + \frac1{y_n}
\]
bildet.
$(y_n)_{n\in\mathbb{N}}$ ist eine Cauchy-Folge, die gegen das Verhältnis
$\varphi=(\!\sqrt{5}+1)/2$ des goldenen Schnittes konvergiert (siehe
Tabelle~\ref{buch:komplex:zahlen:table:gold}, zweite Spalte).
%
% table-gold.tex
%
% (c) 2026 Prof Dr Andreas Müller
%
\begin{table}
\def\p{\phantom{-}}
\centering
\begin{tabular}{|>{$}r<{$}|>{$}l<{$}|>{$}l<{$}|}
\hline
&\begin{minipage}[t]{3.5cm}\raggedright
\setlength{\abovedisplayskip}{0pt}
\setlength{\belowdisplayskip}{0pt}
\raisebox{3pt}{\strut}Rekursive Approxima- tion
\[
 y_{n+1}=1+\frac{1}{y_n}
\]
für $\varphi$
\end{minipage}&\begin{minipage}[t]{3.5cm}\raggedright
Newton-Approximation für Nullstellen von $f(x)=x^2+1$
\end{minipage}
\\[44pt]
\hline
n & y_n\raisebox{3pt}{\mathstrut} & x_n \\[2pt]
\hline
 1 & 2                 & \p0.508034707497476 \raisebox{4pt}{\mathstrut}
\\
 2 & 1.500000000000000             & -0.730167373438399 \\
 3 & 1.\underline{6}66666666666667 &\p0.319690815934174 \\
 4 & 1.\underline{6}00000000000000 & -1.404165739925735 \\
 5 & 1.\underline{6}25000000000000 & -0.345999548896763 \\
 6 & 1.\underline{61}5384615384615 &\p1.272088814812140 \\
 7 & 1.\underline{61}9047619047619 &\p0.242990090617789 \\
 8 & 1.\underline{61}7647058823529 & -1.936202034966178 \\
 9 & 1.\underline{618}181818181818 & -0.709863503540628 \\
10 & 1.\underline{61}7977528089888 &\p0.349429012666962 \\
11 & 1.\underline{6180}55555555556 & -1.256191291052456 \\
12 & 1.\underline{6180}25751072961 & -0.230067094013908 \\
13 & 1.\underline{61803}7135278515 &\p2.058245522488201 \\
14 & 1.\underline{61803}2786885246 &\p0.786197418015100 \\
15 & 1.\underline{61803}4447821682 & -0.242873870579829 \\
16 & 1.\underline{618033}813400125 &\p1.937244794475075 \\
17 & 1.\underline{61803}4055727554 &\p0.710523884635529 \\
18 & 1.\underline{6180339}63166706 & -0.348444169203709 \\
19 & 1.\underline{6180339}98521804 &\p1.260728028475482 \\
20 & 1.\underline{61803398}5017358 &\p0.233767770871423 \\
21 & 1.\underline{6180339}90175597 & -2.021990939507587 \\
22 & 1.\underline{618033988}205325 & -0.763714440828034 \\
23 & 1.\underline{618033988}957902 &\p0.272837745754085 \\
24 & 1.\underline{618033988}670443 & -1.696172136911838 \\
25 & 1.\underline{6180339887}80243 & -0.553304666781480 \\
26 & 1.\underline{6180339887}38303 &\p0.627008940439556 \\
27 & 1.\underline{6180339887}54322 & -0.483932324938967 \\
28 & 1.\underline{61803398874}8204 &\p0.791236155774208 \\
29 & 1.\underline{6180339887}50541 & -0.236304511027909 \\
30 & 1.\underline{618033988749}648 &\p1.997761646531475 \\[2pt]
\hline
\infty & 1.618033988749895 & \raisebox{4pt}{\mathstrut}\\[2pt]
\hline
\end{tabular}
\caption{Approximation der Zahl $\varphi=(1+\!\sqrt{5})/2$
durch die Folge $(y_n)_{n\in\mathbb{N}}$, die durch $y_0=1$ und
$y_{n+1}=1+1/y_n$ rekursiv definiert ist (Spalte $y_n$).
In der Spalte $x_n$ sind die Newton-Approximationen für die
Funktion $f(x) = x^2 + 1$, rekursiv definiert durch
$x_{n+1} = x_n -  f(x_n)/f'(x_n)$, dargestellt, die keinerlei
Anzeichen von Konvergenz zeigen.
\label{buch:komplex:zahlen:table:gold}}
\end{table}


Approximationsfolgen für die Nullstellen einer Funktion $f(x)$ können
zum Beispiel durch das Newton-Verfahren gefunden werden.
\index{Newton-Verfahren}%
\index{Nullstelle}%
Dazu wählt man eine Anfangsschätzung $x_0$ und bildet dann die 
Folge
\[
x_{n+1}
=
x_n - \frac{f(x_n)}{f'(x_n)}.
\]
Falls $x_0$ genügend nahe bei einer Nullstelle $\hat{x}$ liegt, wird die
Folge $(x_n)_{n\in\mathbb{N}}$ gegen die Nullstelle konvergieren:
\[
\lim_{n\to\infty} x_n
=
\hat{x}.
\]
Wendet man dies auf die Funktion $f(x)=x^2+1$ an, welche $f(x) \ge 1$
für alle $x\in\mathbb{R}$ keine reelle Nullstelle hat, erhält man eine
Folge, die keine Cauchy-Folge ist
(Abbildung~\ref{buch:komplex:zahlen:table:gold}, dritte Spalte).
Die Gleichung
\begin{equation}
x^2 + 1 = 0
\qquad\Leftrightarrow\qquad
x^2 = -1
\end{equation}
hat also keine Lösung, die mit einer Approximationsfolge gefunden werden
kann.

